\IfFileExists{stacks-project.cls}{%
\documentclass{stacks-project}
}{%
\documentclass{amsart}
}
\usepackage{amssymb}

% For dealing with references we use the comment environment
\usepackage{verbatim}
\usepackage{fancyvrb}
\newenvironment{reference}{\comment}{\endcomment}
%\newenvironment{reference}{}{}
\newenvironment{slogan}{\comment}{\endcomment}
\newenvironment{history}{\comment}{\endcomment}

% For source figures retained as assets
\usepackage{graphicx}

% For commutative diagrams we use Xy-pic
\usepackage[all]{xy}

% We use 2cell for 2-commutative diagrams.
\xyoption{2cell}
\UseAllTwocells

% We use multicol for the list of chapters between chapters
\usepackage{multicol}

% This is generally recommended for better output
\usepackage{lmodern}
\usepackage[T1]{fontenc}
\usepackage{textalpha}

% For cross-file-references
% We prefer the snapshot tmp/xref/#2.aux,
% which latexmkrc replaces atomically, so
% that documents compiled simultaneously
% never read a half-written .aux file.
% The last argument keeps links to #2.pdf.
\usepackage{xr-hyper}
\newcommand{\maybeexternaldocument}[2]{%
  \IfFileExists{#2.aux}{%
    \externaldocument[#1]{#2}%
  }{%
    \PackageWarning{stack}{Missing #2.aux; external references with prefix #1 unavailable}%
  }%
}
\let\externaldocumentorig\externaldocument
\renewcommand{\externaldocument}[2][]{%
  \IfFileExists{tmp/xref/#2.aux}{%
    \externaldocumentorig[#1]{tmp/xref/#2}[#2.pdf]%
  }{%
  \IfFileExists{#2.aux}{%
    \externaldocumentorig[#1]{#2}%
  }{%
    \PackageWarning{stack}{Missing #2.aux; external references with prefix #1 unavailable}%
  }%
  }%
}

% Package for hypertext links:
\usepackage{hyperref}

% For any local file, say "hello.tex" you want to link to please
% use \externaldocument[hello-]{hello}
\externaldocument[introduction-]{introduction}
\externaldocument[conventions-]{conventions}
\externaldocument[sets-]{sets}
\externaldocument[categories-]{categories}
\externaldocument[topology-]{topology}
\externaldocument[sheaves-]{sheaves}
\externaldocument[sites-]{sites}
\externaldocument[stacks-]{stacks}
\externaldocument[fields-]{fields}
\externaldocument[algebra-]{algebra}
\externaldocument[brauer-]{brauer}
\externaldocument[homology-]{homology}
\externaldocument[derived-]{derived}
\externaldocument[simplicial-]{simplicial}
\externaldocument[more-algebra-]{more-algebra}
\externaldocument[smoothing-]{smoothing}
\externaldocument[modules-]{modules}
\externaldocument[sites-modules-]{sites-modules}
\externaldocument[injectives-]{injectives}
\externaldocument[cohomology-]{cohomology}
\externaldocument[sites-cohomology-]{sites-cohomology}
\externaldocument[dga-]{dga}
\externaldocument[dpa-]{dpa}
\externaldocument[sdga-]{sdga}
\externaldocument[hypercovering-]{hypercovering}
\externaldocument[schemes-]{schemes}
\externaldocument[constructions-]{constructions}
\externaldocument[properties-]{properties}
\externaldocument[morphisms-]{morphisms}
\externaldocument[coherent-]{coherent}
\externaldocument[divisors-]{divisors}
\externaldocument[limits-]{limits}
\externaldocument[varieties-]{varieties}
\externaldocument[topologies-]{topologies}
\externaldocument[descent-]{descent}
\externaldocument[perfect-]{perfect}
\externaldocument[more-morphisms-]{more-morphisms}
\externaldocument[flat-]{flat}
\externaldocument[groupoids-]{groupoids}
\externaldocument[more-groupoids-]{more-groupoids}
\externaldocument[etale-]{etale}
\externaldocument[chow-]{chow}
\externaldocument[intersection-]{intersection}
\externaldocument[pic-]{pic}
\externaldocument[weil-]{weil}
\externaldocument[adequate-]{adequate}
\externaldocument[dualizing-]{dualizing}
\externaldocument[duality-]{duality}
\externaldocument[discriminant-]{discriminant}
\externaldocument[derham-]{derham}
\externaldocument[local-cohomology-]{local-cohomology}
\externaldocument[algebraization-]{algebraization}
\externaldocument[curves-]{curves}
\externaldocument[resolve-]{resolve}
\externaldocument[models-]{models}
\externaldocument[functors-]{functors}
\externaldocument[equiv-]{equiv}
\externaldocument[pione-]{pione}
\externaldocument[etale-cohomology-]{etale-cohomology}
\externaldocument[proetale-]{proetale}
\externaldocument[relative-cycles-]{relative-cycles}
\externaldocument[more-etale-]{more-etale}
\externaldocument[trace-]{trace}
\externaldocument[crystalline-]{crystalline}
\externaldocument[spaces-]{spaces}
\externaldocument[spaces-properties-]{spaces-properties}
\externaldocument[spaces-morphisms-]{spaces-morphisms}
\externaldocument[decent-spaces-]{decent-spaces}
\externaldocument[spaces-cohomology-]{spaces-cohomology}
\externaldocument[spaces-limits-]{spaces-limits}
\externaldocument[spaces-divisors-]{spaces-divisors}
\externaldocument[spaces-over-fields-]{spaces-over-fields}
\externaldocument[spaces-topologies-]{spaces-topologies}
\externaldocument[spaces-descent-]{spaces-descent}
\externaldocument[spaces-perfect-]{spaces-perfect}
\externaldocument[spaces-more-morphisms-]{spaces-more-morphisms}
\externaldocument[spaces-flat-]{spaces-flat}
\externaldocument[spaces-groupoids-]{spaces-groupoids}
\externaldocument[spaces-more-groupoids-]{spaces-more-groupoids}
\externaldocument[bootstrap-]{bootstrap}
\externaldocument[spaces-pushouts-]{spaces-pushouts}
\externaldocument[spaces-chow-]{spaces-chow}
\externaldocument[groupoids-quotients-]{groupoids-quotients}
\externaldocument[spaces-more-cohomology-]{spaces-more-cohomology}
\externaldocument[spaces-simplicial-]{spaces-simplicial}
\externaldocument[spaces-duality-]{spaces-duality}
\externaldocument[formal-spaces-]{formal-spaces}
\externaldocument[restricted-]{restricted}
\externaldocument[spaces-resolve-]{spaces-resolve}
\externaldocument[formal-defos-]{formal-defos}
\externaldocument[defos-]{defos}
\externaldocument[cotangent-]{cotangent}
\externaldocument[examples-defos-]{examples-defos}
\externaldocument[algebraic-]{algebraic}
\externaldocument[examples-stacks-]{examples-stacks}
\externaldocument[stacks-sheaves-]{stacks-sheaves}
\externaldocument[criteria-]{criteria}
\externaldocument[artin-]{artin}
\externaldocument[quot-]{quot}
\externaldocument[stacks-properties-]{stacks-properties}
\externaldocument[stacks-morphisms-]{stacks-morphisms}
\externaldocument[stacks-limits-]{stacks-limits}
\externaldocument[stacks-cohomology-]{stacks-cohomology}
\externaldocument[stacks-perfect-]{stacks-perfect}
\externaldocument[stacks-introduction-]{stacks-introduction}
\externaldocument[stacks-more-morphisms-]{stacks-more-morphisms}
\externaldocument[stacks-geometry-]{stacks-geometry}
\externaldocument[moduli-]{moduli}
\externaldocument[moduli-curves-]{moduli-curves}
\externaldocument[examples-]{examples}
\externaldocument[exercises-]{exercises}
\externaldocument[guide-]{guide}
\externaldocument[desirables-]{desirables}
\externaldocument[coding-]{coding}
\externaldocument[obsolete-]{obsolete}
\externaldocument[fdl-]{fdl}
\externaldocument[index-]{index}

% Heap Project documents
\externaldocument[k3-]{k3}
\externaldocument[torelli-]{torelli}
\externaldocument[linear-algebra-]{linear-algebra}
\externaldocument[groups-]{groups}
\externaldocument[complex-analysis-]{complex-analysis}
\externaldocument[functional-analysis-]{functional-analysis}
\externaldocument[categories-]{categories}
\externaldocument[commutative-algebra-]{commutative-algebra}
\externaldocument[ringed-spaces-]{ringed-spaces}
\externaldocument[homological-algebra-]{homological-algebra}
\externaldocument[lie-algebras-]{lie-algebras}
\externaldocument[representation-theory-]{representation-theory}
\externaldocument[smooth-manifolds-]{smooth-manifolds}
\externaldocument[differential-topology-]{differential-topology}
\externaldocument[algebraic-topology-]{algebraic-topology}
\externaldocument[differential-geometry-]{differential-geometry}
\externaldocument[algebraic-geometry-]{algebraic-geometry}
\externaldocument[symplectic-geometry-]{symplectic-geometry}
\externaldocument[complex-geometry-]{complex-geometry}
\externaldocument[hodge-theory-]{hodge-theory}
\externaldocument[%
variations-of-hodge-structures-]{%
variations-of-hodge-structures}
\externaldocument[%
homological-mirror-symmetry-]{%
homological-mirror-symmetry}
\maybeexternaldocument{deformations-}{deformations}
\externaldocument[vertex-algebras-]{vertex-algebras}
\externaldocument[springer-]{springer}
\externaldocument[infty-categories-]{infty-categories}
\externaldocument[seminars-differential-geometry-]{%
seminars-differential-geometry}
\externaldocument[seminars-complex-geometry-]{%
seminars-complex-geometry}
\externaldocument[seminars-algebraic-geometry-]{%
seminars-algebraic-geometry}
\externaldocument[seminars-others-]{%
seminars-others}
\externaldocument[%
geometric-prequantization-]{%
geometric-prequantization}
\externaldocument[surfaces-]{surfaces}
\externaldocument[birational-maps-commutative-algebra-]{birational-maps-commutative-algebra}
\externaldocument[cremona-transformations-]{cremona-transformations}
\externaldocument[derived-categories-of-sheaves-]{derived-categories-of-sheaves}
\externaldocument[geometric-invariant-theory-]{geometric-invariant-theory}
\externaldocument[geometric-stability-conditions-and-group-actions-]{geometric-stability-conditions-and-group-actions}
\externaldocument[moduli-spaces-of-sheaves-]{moduli-spaces-of-sheaves}
\externaldocument[bridgeland-stability-general-theory-]{bridgeland-stability-general-theory}
\externaldocument[hilbert-schemes-]{hilbert-schemes}
\externaldocument[noncommutative-abelian-surfaces-and-kummer-type-hyperkahler-varieties-]{noncommutative-abelian-surfaces-and-kummer-type-hyperkahler-varieties}
\externaldocument[mumford-tate-groups-in-hodge-theory-]{mumford-tate-groups-in-hodge-theory}
\externaldocument[hodge-birational-atoms-2026-]{hodge-birational-atoms-2026}
\externaldocument[xxii-egd-2026-]{%
xxii-egd-2026}
\externaldocument[pde-]{pde}
\externaldocument[probability-]{probability}

\externaldocument[linguistics-]{linguistics}
\externaldocument[genealogia-familiar-]{%
genealogia-familiar}

\externaldocument[%
16th-alga-meeting-2026-]{%
16th-alga-meeting-2026}

\externaldocument[%
iv-jornadas-lefschetz-2026-]{%
iv-jornadas-lefschetz-2026}

\externaldocument[reading-the-masters-]{%
reading-the-masters}

\externaldocument[real-analysis-]{%
real-analysis}

\externaldocument[optimal-transport-]{%
optimal-transport}

\externaldocument[history-]{history}

\externaldocument[seiberg-witten-]{seiberg-witten}

\externaldocument[history-philosophy-biology-]{history-philosophy-biology}

\externaldocument[chemistry-]{chemistry}

\externaldocument[molecular-cell-biology-]{molecular-cell-biology}

\externaldocument[microbiology-]{microbiology}

\externaldocument[protistology-]{protistology}

\externaldocument[mycology-]{mycology}

\externaldocument[botany-]{botany}

\externaldocument[zoology-]{zoology}

\externaldocument[animal-histology-]{animal-histology}

\externaldocument[animal-physiology-]{animal-physiology}

\externaldocument[systematics-biogeography-]{systematics-biogeography}

\externaldocument[ecology-natural-resources-]{ecology-natural-resources}

\externaldocument[earth-sciences-]{earth-sciences}

\externaldocument[biostatistics-]{biostatistics}

\externaldocument[mathematical-systems-biology-]{mathematical-systems-biology}

\externaldocument[cimpa-el-salvador-2024-]{cimpa-el-salvador-2024}

% Theorem environments.
%
\theoremstyle{plain}
\newtheorem{theorem}[subsection]{Theorem}
\newtheorem{proposition}[subsection]{Proposition}
\newtheorem{lemma}[subsection]{Lemma}
\newtheorem{corollary}[subsection]{Corollary}

\theoremstyle{definition}
\newtheorem{definition}[subsection]{Definition}
\newtheorem{example}[subsection]{Example}
\newtheorem{exercise}[subsection]{Exercise}
\newtheorem{situation}[subsection]{Situation}

\theoremstyle{remark}
\newtheorem{remark}[subsection]{Remark}
\newtheorem{remarks}[subsection]{Remarks}

\numberwithin{equation}{subsection}

% Macros
%
\def\lim{\mathop{\mathrm{lim}}\nolimits}
\def\colim{\mathop{\mathrm{colim}}\nolimits}
\def\Spec{\mathop{\mathrm{Spec}}}
\def\Hom{\mathop{\mathrm{Hom}}\nolimits}
\def\Ext{\mathop{\mathrm{Ext}}\nolimits}
\def\SheafHom{\mathop{\mathcal{H}\!\mathit{om}}\nolimits}
\def\SheafExt{\mathop{\mathcal{E}\!\mathit{xt}}\nolimits}
\def\Sch{\mathit{Sch}}
\def\Mor{\mathop{\mathrm{Mor}}\nolimits}
\def\Ob{\mathop{\mathrm{Ob}}\nolimits}
\def\Sh{\mathop{\mathit{Sh}}\nolimits}
\def\NL{\mathop{N\!L}\nolimits}
\def\CH{\mathop{\mathrm{CH}}\nolimits}
\def\proetale{{pro\text{-}\acute{e}tale}}
\def\etale{{\acute{e}tale}}
\def\QCoh{\mathit{QCoh}}
\def\Ker{\mathop{\mathrm{Ker}}}
\def\Im{\mathop{\mathrm{Im}}}
\def\Coker{\mathop{\mathrm{Coker}}}
\def\Coim{\mathop{\mathrm{Coim}}}

% Boxtimes
%
\DeclareMathSymbol{\boxtimes}{\mathbin}{AMSa}{"02}

%
% Macros for moduli stacks/spaces
%
\def\QCohstack{\mathcal{QC}\!\mathit{oh}}
\def\Cohstack{\mathcal{C}\!\mathit{oh}}
\def\Spacesstack{\mathcal{S}\!\mathit{paces}}
\def\Quotfunctor{\mathrm{Quot}}
\def\Hilbfunctor{\mathrm{Hilb}}
\def\Curvesstack{\mathcal{C}\!\mathit{urves}}
\def\Polarizedstack{\mathcal{P}\!\mathit{olarized}}
\def\Complexesstack{\mathcal{C}\!\mathit{omplexes}}
% \Pic is the operator that assigns to X its picard group, usage \Pic(X)
% \Picardstack_{X/B} denotes the Picard stack of X over B
% \Picardfunctor_{X/B} denotes the Picard functor of X over B
\def\Pic{\mathop{\mathrm{Pic}}\nolimits}
\def\Picardstack{\mathcal{P}\!\mathit{ic}}
\def\Picardfunctor{\mathrm{Pic}}
\def\Deformationcategory{\mathcal{D}\!\mathit{ef}}


\begin{document}

\title{Zoology}
\maketitle

\phantomsection
\label{section-phantom}

\tableofcontents

\section{What this chapter covers}
\label{section-overview}

These notes survey the animal kingdom as
Daniel studied it: the general architecture
of animal bodies, the invertebrate phyla from
the mesozoans to the arthropods, molluscs and
annelids, and the diversity of several of
these groups in Mexico. Two short sections,
on embryology and on the vertebrates,
complete the picture from the textbooks.

\textbf{Sources of these notes.} The notes
rest on the course \emph{Animales I} (fourth
semester, Facultad de Ciencias, UNAM, spring
2015), a survey of the invertebrates taught
in three blocks, each closed by an exam. What
survives is Daniel's own work: notes on
animal architecture, three exam outlines (one
per block), slide decks on the
blastocoelomates, the arthropod subphyla, the
paraarthropods (Onychophora and Tardigrada)
and the molluscs, eleven one-page essays
summarizing articles on Mexican biodiversity
from the 2014 supplement of the \emph{Revista
Mexicana de Biodiversidad}, and a list of the
animals seen at an aquarium, from sharks and
rays to the upside-down jellyfish
\emph{Cassiopea xamachana}, the tropical gar
\emph{Atractosteus tropicus} and the axolotl.
The facts were checked and completed with
Hickman, Roberts and Larson, \emph{Animal
Diversity} \cite{hickman2002animal}, Brusca
and Brusca, \emph{Invertebrates}
\cite{brusca2003invertebrates}, Hill, Wyse
and Anderson for osmoregulation
\cite{hill2012animal}, and the biodiversity
articles themselves.

\section{Animal architecture}
\label{section-architecture}

A body plan can be described by a few broad
features: symmetry, the number of germ
layers, the presence and kind of body cavity,
segmentation and cephalization \cite[Chapter
3]{hickman2002animal}.

\textbf{Symmetry.} \emph{Spherical} symmetry
(any plane through the centre gives mirror
halves) is found chiefly in protozoans.
\emph{Radial} symmetry (more than two planes
through the oral--aboral axis) characterizes
the Radiata, Cnidaria and Ctenophora;
ctenophores are strictly \emph{biradial},
since their pair of tentacles leaves only two
planes of symmetry. Echinoderms are
secondarily radial as adults; their larvae
are bilateral. \emph{Bilateral} symmetry (one
sagittal plane) defines the Bilateria and
goes with directional movement and
\emph{cephalization}, the concentration of
nervous tissue and sense organs in a head.
Positions are named anterior and posterior,
dorsal and ventral, medial and lateral; the
sagittal, frontal and transverse planes cut
the body into right and left, dorsal and
ventral, and anterior and posterior parts.

\begin{definition}
\label{definition-germ-layers}
An animal is \emph{diploblastic} if its body
develops from two germ layers, ectoderm and
endoderm, and \emph{triploblastic} if a third
layer, mesoderm, forms between them.
\end{definition}

Cnidarians and ctenophores are diploblastic;
the gelatinous mesoglea between their two
epithelia is not a germ layer. All Bilateria
are triploblastic.

\begin{definition}
\label{definition-coelom}
A \emph{coelom} is a fluid-filled body cavity
between the body wall and the gut that arises
within the mesoderm and is lined by a
mesodermal \emph{peritoneum}. A triploblastic
animal is \emph{acoelomate} if the space
between epidermis and gut is filled with
mesodermal parenchyma, \emph{pseudocoelomate}
(or \emph{blastocoelomate}) if the cavity is
a persistent blastocoel without a peritoneum,
and \emph{eucoelomate} if it has a true
coelom.
\end{definition}

Flatworms are acoelomate; nematodes, rotifers
and their allies are blastocoelomate;
molluscs, annelids, arthropods, echinoderms
and chordates are eucoelomate. A body cavity
gives a tube-within-a-tube body, room for the
viscera, and in many worms a hydrostatic
skeleton \cite[Chapter 3]{hickman2002animal}.

\textbf{Segmentation.} \emph{Metamerism} is
the serial repetition of body units
(metameres or somites), external and
internal, along the anteroposterior axis.
True metamerism occurs in annelids,
arthropods and chordates; it is homologous in
the first two, while chordates evolved it
independently \cite[Chapter
3]{hickman2002animal}.

\section{Organ systems in outline}
\label{section-functions}

This is the content of Daniel's notes on
animal architecture, which list the
functional options open to an animal body.

\textbf{Movement and support.} Amoeboid
movement uses pseudopodia, with a stiff outer
ectoplasm and a flowing inner endoplasm;
small animals and larvae glide on cilia and
flagella. Skeletons are hydrostatic (muscles
acting on an enclosed fluid, as in worms) or
rigid, either an exoskeleton (arthropods) or
an endoskeleton (echinoderms, vertebrates).

\textbf{Feeding.} Digestion may be
extracorporeal (enzymes released onto the
food, as in spiders), extracellular (in a
gut) or intracellular (in food vacuoles, as
in sponges). Feeding strategies include
suspension feeding (trapping particles with
cilia, tentacles or mucus), deposit feeding
(eating sediment and its organic matter),
herbivory, predation and scavenging, uptake
of dissolved organic matter through the body
surface, and reliance on chemoautotrophic
symbionts.

\textbf{Circulation and gas exchange.} In a
\emph{closed} circulatory system the blood
stays in vessels and reaches the tissues
through capillaries (annelids, cephalopods,
vertebrates); in an \emph{open} system it
leaves the vessels and bathes the organs in a
hemocoel (arthropods, most molluscs). Pumps
range from contractile vessels to ostiate
hearts (arthropods, which take blood back
from the hemocoel through ostia) and
chambered hearts. Gas exchange occurs across
gills, tracheae (insects, myriapods,
onychophorans), book lungs (arachnids) or the
vascularized mantle cavity of pulmonate
snails. The respiratory pigments are
hemoglobins, hemocyanins (molluscs and
arthropods), hemerythrins and chlorocruorins
(some annelids).

\textbf{Excretion and osmoregulation.} The
nitrogenous wastes are ammonia (most aquatic
animals), urea and uric acid. An animal that
is hyperosmotic to its medium, as in fresh
water, tends to gain water by osmosis and
lose ions by diffusion; one that is
hyposmotic, as a marine teleost, tends to
lose water and gain ions. Most marine
invertebrates are essentially isosmotic to
sea water and face no osmotic problem
\cite[Chapter 28]{hill2012animal}. The organs
involved are contractile vacuoles in protists
and freshwater sponges, protonephridia (flame
cells) in flatworms, rotifers and nemerteans,
metanephridia in annelids and molluscs,
Malpighian tubules in insects, myriapods and
arachnids, and antennal or maxillary glands
in crustaceans.

\textbf{Nervous system and receptors.} A
stimulus reaches a receptor, travels along
afferent fibres to an integrating centre and
returns along efferent fibres to an effector.
Nervous systems range from the nerve net of
cnidarians and ctenophores to ganglionated
cords and a brain. Receptors are tactile,
georeceptors (statocysts), proprioceptors,
phonoreceptors (tympanal organs),
baroreceptors, chemoreceptors, photoreceptors
and thermoreceptors. Hormones coordinate the
body of one individual; pheromones carry
signals between individuals.

\textbf{Reproduction.} Asexual reproduction
(budding, fission, gemmules) is important in
sponges, cnidarians and many small worms;
regeneration of a lost limb is not
reproduction. Sexual reproduction recombines
genomes. \emph{Parthenogenesis}, development
of an unfertilized egg, is common in rotifers
and in some arthropods.

\section{Animal embryology}
\label{section-embryology}

This section is written from the textbooks;
the surviving coursework has no material on
it.

\textbf{Eggs and cleavage.} Eggs are
polarized along an animal--vegetal axis and
classified by their yolk: \emph{isolecithal}
(little, evenly spread), \emph{telolecithal}
(concentrated at the vegetal pole) or
\emph{centrolecithal} (in the centre, as in
arthropods). Cleavage, the first divisions of
the zygote into blastomeres, is
\emph{holoblastic} (complete) in isolecithal
and moderately telolecithal eggs and
\emph{meroblastic} (incomplete) when much
yolk is present \cite[Chapter
4]{brusca2003invertebrates}. The two classic
patterns in the Bilateria are radial
cleavage, with tiers of cells stacked over
one another, and spiral cleavage, with the
spindles oblique to the axis so that each
quartet of cells lies in the furrows of the
one below \cite[Chapter
3]{hickman2002animal}.

\textbf{Blastula and gastrulation.} Cleavage
produces a blastula: a hollow
\emph{coeloblastula} with a blastocoel, a
solid \emph{stereoblastula}, or a
\emph{discoblastula}, a cap of cells on an
uncleaved yolk mass. \emph{Gastrulation}
forms the germ layers. Its main modes are
invagination (a pouch pushes into the
blastocoel and becomes the archenteron,
opening by the blastopore), ingression (cells
migrate into the blastocoel, as in many
cnidarians), delamination (cells divide
parallel to the surface), epiboly (animal
cells overgrow the vegetal ones) and
involution (cells at the edge of a disc roll
under it). Ectoderm forms the epidermis and
nervous system, endoderm the gut lining, and
mesoderm the muscles, the circulatory system,
most supporting tissue and the coelomic
lining \cite[Chapter
4]{brusca2003invertebrates}.

\textbf{Protostomes and deuterostomes.} The
Bilateria divide into two lineages by a set
of developmental tendencies \cite[Chapter
3]{hickman2002animal}:

\begin{center}
\begin{tabular}{lll}
& Protostomes & Deuterostomes \\
\hline
Cleavage & spiral & radial \\
Development & mosaic & regulative \\
Blastopore & mouth & anus \\
Coelom & schizocoelous & enterocoelous \\
Examples & annelids, & echinoderms, \\
& molluscs, & hemichordates, \\
& arthropods & chordates \\
\end{tabular}
\end{center}

In \emph{mosaic} development the fates of the
blastomeres are fixed early, and an isolated
blastomere gives a partial embryo; in
\emph{regulative} development an isolated
early blastomere can still give a whole,
smaller embryo. In spiral cleavers a single
cell, the 4d mesentoblast, gives all the
mesoderm \cite[Chapter 3]{hickman2002animal}.

\textbf{Coelom formation.} In
\emph{schizocoely} the mesoderm grows as
solid bands from the mesentoblast, which then
split to open the coelom; in
\emph{enterocoely} the mesoderm and coelom
arise together as pouches pinched off the
archenteron, often giving three pairs of
cavities (protocoel, mesocoel, metacoel)
\cite[Chapter 4]{brusca2003invertebrates}.
The final coeloms are indistinguishable. The
vertebrates modify all these tendencies:
their cleavage is rotational in mammals and
discoidal in reptiles, birds and many fishes,
and their coelom forms by a derived,
essentially schizocoelous route \cite[Chapter
3]{hickman2002animal}.

\section{Mesozoans, placozoans and sponges}
\label{section-early-metazoa}

The first block of the course began with
animals that lack true tissues and organs and
do not gastrulate.

\textbf{Mesozoans.} The name Mesozoa is now
used for a grade rather than a phylum: Brusca
and Brusca split it into the Rhombozoa and
the Orthonectida \cite[Chapter
7]{brusca2003invertebrates}. Rhombozoans
(dicyemids) are tiny symbionts in the kidneys
of cephalopods: an outer layer of ciliated
somatic cells, constant in number, surrounds
a single long axial cell inside which new
individuals develop; the vermiform adult
(nematogen) attaches to the host tissue by
its anterior polar cap. Orthonectids are
parasites of invertebrates such as brittle
stars, bivalves, nemerteans, flatworms and
polychaetes. Their asexual phase is an
amoeboid multinucleate plasmodium in the
host's tissues; it produces free-swimming
ciliated males and females.

\textbf{Placozoa.} \emph{Trichoplax
adhaerens} is a marine plate of a few
thousand cells, 2--3 mm across, without
symmetry or anterior--posterior polarity. A
dorsal layer of flat ciliated cells and a
ventral layer of columnar cells enclose a
layer of amoeboid cells in a gel; the ventral
layer feeds by phagocytosis of detritus. It
reproduces by fission and by budding of
swarmers, and sexual stages with eggs are
known \cite[Chapter
7]{brusca2003invertebrates}.

\textbf{Porifera.} Sponges are sessile filter
feeders built around a system of water canals
\cite[Chapter 6]{hickman2002animal}. Water
enters through pores (ostia), is driven by
the flagellated \emph{choanocytes} (collar
cells), and leaves by an osculum. The body
has an outer \emph{pinacoderm} of flat
pinacocytes, with porocytes forming the
pores, an inner \emph{choanoderm}, and
between them the \emph{mesohyl}, a protein
matrix with amoebocytes, totipotent
\emph{archaeocytes} and the skeleton of
spicules (calcium carbonate or silica) and
collagen (spongin). There are no organs or
true tissues, digestion is intracellular, and
the larvae swim. Freshwater sponges survive
bad seasons as \emph{gemmules}, coats of
spicules filled with archaeocytes.

\begin{center}
\begin{tabular}{ll}
Canal system & Organization \\
\hline
Asconoid & tube lined by \\
& choanocytes \\
Syconoid & choanocytes in \\
& radial canals \\
Leuconoid & choanocytes in \\
& flagellated chambers \\
\end{tabular}
\end{center}

Asconoid and syconoid sponges occur only in
the Calcarea; most sponges are leuconoid, the
plan that allows large size. The classes are
Calcarea (calcium carbonate spicules),
Hexactinellida (glass sponges, six-rayed
siliceous spicules) and Demospongiae
(siliceous spicules and spongin, all
leuconoid), which contains most species
\cite[Chapter 6]{hickman2002animal}.

\section{Cnidarians and ctenophores}
\label{section-radiata}

\textbf{Cnidaria.} Cnidarians are radial and
diploblastic: an outer epidermis and an inner
gastrodermis are separated by the mesoglea.
They catch prey with \emph{cnidocytes}, cells
that discharge a stinging capsule
(nematocyst). The gut is a gastrovascular
cavity with a single opening, and the nervous
system is a nerve net. Many alternate between
an asexual, sessile polyp and a sexual,
swimming medusa \cite[Chapter
7]{hickman2002animal}.

\begin{center}
\begin{tabular}{ll}
Class & Diagnostic features \\
\hline
Hydrozoa & polyp usually dominant, \\
& often colonial with \\
& feeding, reproductive \\
& and defensive polyps; \\
& medusa with a velum \\
Scyphozoa & medusa dominant, \\
& no velum, rhopalia; \\
& polyp strobilates \\
& into ephyrae \\
Cubozoa & box-shaped medusa \\
& with a velarium, \\
& rhopalia with eyes, \\
& tentacles on pedalia \\
Anthozoa & polyps only; pharynx \\
& and mesenteries, \\
& often in sixes or \\
& eights; corals with \\
& calcareous skeleton \\
\end{tabular}
\end{center}

Hydrozoa, Scyphozoa and Cubozoa together form
the Medusozoa, the cnidarians with a medusa
stage \cite{gasca2014medusozoa}.

\textbf{Ctenophora.} Comb jellies are marine,
fewer than a hundred species, and swim with
eight rows of ciliary comb plates. They lack
cnidocytes (with one exception that steals
them from its prey) and capture prey with
\emph{colloblasts}, adhesive cells on the
tentacles. An aboral statocyst controls the
comb rows. All are monoecious and there is no
alternation of generations \cite[Chapter
7]{hickman2002animal}.

\section{Flatworms and ribbon worms}
\label{section-acoelomates}

\textbf{Platyhelminthes.} Flatworms are
bilateral, acoelomate and dorsoventrally
flattened, with a gut that has a mouth but no
anus, protonephridia, a ganglionated nervous
system and, in most species, hermaphroditic
reproduction. The parasitic groups are
covered by a syncytial tegument \cite[Chapter
8]{hickman2002animal},
\cite{garciaprieto2014platyhelminthes}.

\begin{center}
\begin{tabular}{ll}
Group & Diagnostic features \\
\hline
Turbellaria & free-living, ciliated \\
& epidermis, eyespots \\
Monogenoidea & ectoparasites (often \\
& on fish gills), one \\
& host, posterior \\
& attachment organ \\
& (haptor) with hooks \\
Trematoda & flukes: oral sucker \\
& and ventral sucker \\
& (acetabulum); snail \\
& as first host \\
Cestoda & tapeworms: scolex, \\
& neck and strobila \\
& of proglottids; no gut \\
\end{tabular}
\end{center}

\textbf{Nemertea.} Ribbon worms share with
flatworms a ciliated epidermis and flame
cells but differ in having a complete gut
with an anus, a closed circulatory system,
and an eversible proboscis lying in its own
cavity, the rhynchocoel, above the gut. Most
are marine carnivores and the sexes are
usually separate \cite[Chapter
8]{hickman2002animal}.

\section{Blastocoelomates}
\label{section-blastocoelomates}

The second block of the course dealt with
small worms whose body cavity is a persistent
blastocoel. The slide deck lists their shared
features: bilateral symmetry, no
segmentation, a complete gut with a well
developed pharynx, a body wall of cellular or
syncytial epidermis under a cuticle that is
often moulted, no respiratory or circulatory
organs, protonephridia in many, and frequent
\emph{eutely}, a fixed number of cells or
nuclei in the adult. The group is a grade,
not a clade: sequence data assign nematodes,
nematomorphs and kinorhynchs to the moulting
Ecdysozoa, with the arthropods, and rotifers,
acanthocephalans, gastrotrichs and entoprocts
to the Lophotrochozoa, with molluscs and
annelids; the place of loriciferans is
undetermined \cite[Chapter
9]{hickman2002animal}.

\begin{center}
\begin{tabular}{ll}
Phylum & Diagnostic features \\
\hline
Rotifera & ciliated corona; \\
& mastax with trophi; \\
& adhesive toes; \\
& parthenogenesis \\
Gastrotricha & adhesive tubes; \\
& cuticular scales or \\
& spines; hermaphrodites \\
Kinorhyncha & 13 segments with \\
& cuticular plates; \\
& spines; moulting \\
Nematoda & cylindrical; cuticle; \\
& longitudinal muscles \\
& only; four moults \\
Nematomorpha & larvae parasitic in \\
& arthropods, adults \\
& free-living \\
Acanthocephala & spiny proboscis; \\
& no gut; parasites \\
Entoprocta & sessile; tentacle \\
& crown around mouth \\
& and anus; U-shaped gut \\
Loricifera & body retracts into \\
& an abdominal lorica; \\
& Higgins larva \\
Cycliophora & \emph{Symbion} on \\
& lobster mouthparts; \\
& dwarf males \\
\end{tabular}
\end{center}

\textbf{Nematoda.} The cross-section Daniel
drew for the exam shows a cuticle, a
syncytial epidermis thickened into dorsal,
ventral and lateral cords, a single layer of
longitudinal muscle cells whose processes
reach the dorsal and ventral nerve cords, and
the gut and gonads in the fluid-filled
blastocoel. A nerve ring surrounds the
pharynx. Without circular muscles, nematodes
move by undulation against the pressure of
the blastocoel fluid. Males have copulatory
spicules, and there are four juvenile stages;
in parasitic species the third-stage larva is
usually the one that infects the definitive
host \cite[Chapter 9]{hickman2002animal},
\cite{garciaprieto2014nematoda}.

\textbf{Acanthocephala.} All are intestinal
parasites of vertebrates as adults, with
arthropods as intermediate hosts. They attach
by a proboscis armed with hooks and absorb
food through the body wall; the males have
cement glands that plug the female after
copulation
\cite{garciaprieto2014acanthocephala}.

\textbf{Rotifera.} About 1800 species, mostly
in fresh water. The anterior \emph{corona} of
cilia creates the illusion of turning wheels;
the muscular pharynx (mastax) carries hard
jaws (trophi). The epidermis is syncytial
with a fixed number of nuclei. The classes
are Seisonidea, Bdelloidea (only
parthenogenetic females known) and
Monogononta (males reduced and appearing
seasonally) \cite[Chapter
12]{brusca2003invertebrates}.

\section{Arthropods and their allies}
\label{section-arthropods}

The third block of the course dealt with the
arthropods, molluscs and annelids.

\textbf{Arthropoda.} Arthropods have a
chitinous exoskeleton, calcified in
crustaceans and millipedes, which is moulted
(ecdysis); segments fused into functional
regions (\emph{tagmosis}); jointed
appendages; no cilia; a hemocoel with a
dorsal ostiate heart; a complete gut; a
ventral ganglionated nerve cord; and gas
exchange by gills, tracheae or book lungs.
Daniel's outline lists the layers of the
crustacean cuticle: epicuticle, exocuticle,
calcified and uncalcified endocuticle, over
the epidermis and its basement membrane
\cite[Chapter 12]{hickman2002animal}.

\begin{center}
\begin{tabular}{ll}
Subphylum & Diagnostic features \\
\hline
Cheliceriformes & prosoma and \\
& opisthosoma; \\
& chelicerae; \\
& no antennae \\
Myriapoda & head and long \\
& trunk; one pair of \\
& antennae; mandibles \\
Crustacea & two pairs of \\
& antennae; biramous \\
& limbs; nauplius \\
Hexapoda & head, thorax, \\
& abdomen; three \\
& pairs of legs \\
\end{tabular}
\end{center}

\textbf{Chelicerates.} The prosoma carries
six pairs of appendages: chelicerae,
pedipalps and four pairs of legs; the
opisthosoma usually lacks appendages.
Pycnogonida (sea spiders) are marine
carnivores with a proboscis and a reduced
opisthosoma; Merostomata (horseshoe crabs)
have a rigid prosomal shield, book gills on
the opisthosoma and a telson; Arachnida are
mostly terrestrial and develop directly. The
arachnid orders seen in the course were
scorpions, spiders, pseudoscorpions,
solifuges, uropygids, schizomids, harvestmen,
palpigrades, ricinuleids and the mites, whose
body is fused into a single region
(idiosoma). Mites are the most diverse
arachnids; ticks are their best known group
in Mexico \cite{perez2014acari}.

\textbf{Myriapods.} Head and trunk, one pair
of antennae, uniramous legs, no compound
eyes, tracheae and Malpighian tubules, direct
development \cite[Chapter
18]{brusca2003invertebrates}. Centipedes
(Chilopoda) are flattened predators whose
first trunk appendages are poison claws
(forcipules); millipedes (Diplopoda) are
cylindrical, with a calcified cuticle and
most trunk segments fused in pairs
(diplosegments), and many roll up; Symphyla
have twelve leg pairs; Pauropoda are minute,
with branched antennae and no tracheae.

\textbf{Crustaceans.} Mandibulate, mostly
aquatic, with two pairs of antennae, biramous
appendages, gills, antennal or maxillary
excretory glands and, in marine forms, a
nauplius larva. The classes in the course
deck were Remipedia, Cephalocarida,
Branchiopoda (\emph{Artemia},
\emph{Daphnia}), Ostracoda (body inside a
bivalved carapace), Mystacocarida, Cirripedia
(barnacles, sessile filter feeders inside
plates), Copepoda, Branchiura (fish lice) and
Malacostraca (eight thoracic and six
abdominal segments; decapods, isopods,
krill). The deck also listed the tongue worms
(Pentastomida), endoparasites of the
respiratory tract of vertebrates; molecular
data show them to be highly modified
crustaceans, perhaps derived from Branchiura
\cite[Chapter 16]{brusca2003invertebrates}.

\textbf{Hexapods.} Head with antennae,
mandibles, maxillae and labium; thorax with
three pairs of legs; abdomen without legs.
The Entognatha (Collembola, Protura, Diplura)
have the mouthparts hidden in the head and
are wingless; the Insecta usually have wings
and breathe by tracheae. By development,
insects are \emph{ametabolous} (young like
small adults, as silverfish),
\emph{hemimetabolous} (gradual metamorphosis
through nymphs whose wings develop
externally, as bugs, grasshoppers and the
aquatic nymphs of mayflies and dragonflies)
or \emph{holometabolous} (larva, pupa and
adult, as beetles, butterflies, flies and
bees) \cite[Chapter 12]{hickman2002animal}.
Beetles (Coleoptera, with hardened forewings,
the elytra) form the largest order of
animals.

\textbf{Onychophora and Tardigrada.}
Onychophorans (velvet worms, about 110
species) have a head with antennae, jaws and
oral papillae, and a trunk of unjointed legs
(lobopods) ending in claws; a soft chitinous
cuticle that is moulted; a hemocoel;
tracheae; a pair of nephridia per segment; a
ladder-like nervous system; separate sexes
and direct development. Tardigrades (water
bears, about 800 species in the course deck)
are microscopic, with four pairs of clawed
lobopods, eutely and the ability to survive
drying by cryptobiosis. Both moult and are
close relatives of the arthropods
\cite[Chapter 13]{hickman2002animal}.

\section{Molluscs and annelids}
\label{section-molluscs-annelids}

\textbf{Mollusca.} The phylum is defined by
five synapomorphies: mantle, foot, radula, a
characteristic nervous system and comb-like
gills (ctenidia) \cite{castillo2014moluscos}.
Molluscs are protostome coelomates with
metanephridia and, except cephalopods, an
open circulation. Eight classes are living,
and two (Rostroconchia, Helcionelloida) are
extinct.

\begin{center}
\begin{tabular}{ll}
Class & Diagnostic features \\
\hline
Solenogastres, & wormlike, no shell, \\
Caudofoveata & cuticle with spicules \\
& (``Aplacophora'') \\
Monoplacophora & one cap-shaped shell; \\
& serially repeated \\
& gills and nephridia \\
Polyplacophora & chitons: eight \\
& dorsal plates \\
Scaphopoda & tusk shell open at \\
& both ends; no gills \\
Gastropoda & torsion of the \\
& visceral mass; \\
& often an operculum \\
Bivalvia & two valves and \\
& a ligament; filter \\
& feeders; no radula \\
Cephalopoda & foot as arms and \\
& funnel; beak and \\
& radula; large brain \\
\end{tabular}
\end{center}

In gastropods, \emph{torsion} rotates the
visceral mass 180 degrees during development,
bringing the mantle cavity forward; many lose
the right gill and nephridium. Cephalopods
comprise Nautiloidea, with an external
chambered shell in which the animal occupies
the last chamber and a siphuncle connects the
others, and Coleoidea (squids, cuttlefish,
octopuses), with an internal or absent shell,
an ink sac and statocysts \cite[Chapter
10]{hickman2002animal}. In male coleoids a
modified arm, the hectocotylus, transfers the
spermatophores \cite[Chapter
20]{brusca2003invertebrates}.

\textbf{Annelida.} Segmented worms have a
metameric body with repeated coelomic
compartments, nephridia, ganglia and blood
vessels; chitinous bristles (setae, or
chaetae); a closed circulation with dorsal
and ventral vessels; a complete, unsegmented
gut; and, in marine forms, a trochophore
larva. The head is formed by the prostomium
and the peristomium \cite[Chapter
11]{hickman2002animal}.

\begin{center}
\begin{tabular}{ll}
Group & Diagnostic features \\
\hline
Polychaeta & parapodia with many \\
& setae; no clitellum; \\
& usually dioecious; \\
& trochophore \\
Oligochaeta & few setae; no \\
& parapodia; clitellum; \\
& hermaphrodites; \\
& direct development \\
Hirudinea & 34 segments; two \\
& suckers; no setae; \\
& clitellum; \\
& hermaphrodites \\
\end{tabular}
\end{center}

The clitellum, a glandular girdle that
secretes the cocoon, unites oligochaetes and
leeches as the clitellates
\cite{oceguera2014euhirudinea},
\cite{fragoso2014oligochaeta}. Molecular
phylogenies have since expanded the Annelida
to include groups once treated as separate
phyla (Echiura, Sipuncula, Pogonophora) and
place the clitellates among the polychaetes;
the traditional three classes are kept here
as in the course
\cite{tovarhernandez2014polychaeta}.

\section{Animal diversity in Mexico}
\label{section-mexico}

Volume 85 of the \emph{Revista Mexicana de
Biodiversidad} (2014) carried a supplement of
inventories of the Mexican biota. Daniel
summarized the articles on sponges,
medusozoans, parasitic flatworms and
nematodes, acanthocephalans, annelids, marine
molluscs, mites, weevils and fishes; the
course also gave out the articles on decapods
and butterflies. Their headline counts:

\begin{center}
\begin{tabular}{lrr}
Group & Species & World \\
\hline
Demospongiae
\cite{carballo2014porifera}
& 517 &  \\
Medusozoa
\cite{gasca2014medusozoa}
& 289 & 7.3\% \\
Parasitic flatworms
\cite{garciaprieto2014platyhelminthes}
& 1015 & 7.4\% \\
Parasitic nematodes
\cite{garciaprieto2014nematoda}
& 402 & 4.8\% \\
Acanthocephala
\cite{garciaprieto2014acanthocephala}
& 60 & 5\% \\
Polychaeta
\cite{tovarhernandez2014polychaeta}
& 1500 &  \\
Earthworms
\cite{fragoso2014oligochaeta}
& 102 &  \\
Leeches
\cite{oceguera2014euhirudinea}
& 31 & 4.5\% \\
Decapoda
\cite{alvarez2014decapoda}
& 1775 & 11.9\% \\
Mites and ticks
\cite{perez2014acari}
& 2625 & 4.8\% \\
Curculionoidea
\cite{morrone2014curculionoidea}
& 3594 &  \\
Marine molluscs
\cite{castillo2014moluscos}
& 4643 &  \\
Fishes
\cite{espinosa2014peces}
& 2763 & 9.8\% \\
\end{tabular}
\end{center}

Some details behind these numbers:

\begin{itemize}
\item The sponge count covers only the
Demospongiae, since Calcarea and
Hexactinellida are almost unstudied in
Mexico; 174 of the 517 species are endemic,
292 occur on the Atlantic side, but most
endemics are Pacific
\cite{carballo2014porifera}.
\item The medusozoans comprise 3 cubozoans,
20 scyphozoans and 266 hydrozoans
\cite{gasca2014medusozoa}.
\item The parasites are those of wild
vertebrates. The flatworms are 544
trematodes, 269 monogeneans and 202 cestodes,
recorded from 1249 host species; Veracruz has
the most (185), and the authors estimate that
more than 3000 species remain to be recorded
\cite{garciaprieto2014platyhelminthes}. Of
the nematodes, the most studied hosts are
ray-finned fishes but the richest fauna is
that of mammals (166 species)
\cite{garciaprieto2014nematoda}. Of the 5488
vertebrate species recorded in Mexico, 94.5\%
have never been examined for acanthocephalans
\cite{garciaprieto2014acanthocephala}.
\item The 1500 polychaetes belong to 63
families, the largest being Nereididae and
Syllidae; about 13\% may be endemic
\cite{tovarhernandez2014polychaeta}. Of the
102 described earthworms, 51 are native and
51 exotic; native species dominate natural
habitats and exotics disturbed ones, and the
authors estimate about 130 native species in
all \cite{fragoso2014oligochaeta}. Twenty of
the 31 leeches are endemic, out of about 680
species worldwide
\cite{oceguera2014euhirudinea}.
\item The decapods are 1597 marine and 178
freshwater species; marine endemism is low,
but freshwater crabs and crayfish are highly
endemic \cite{alvarez2014decapoda}.
\item The mites belong to five orders and are
4.8\% of 54\,617 species worldwide
\cite{perez2014acari}. The weevils, with a
rostrum carrying the mouthparts, are about
62\,000 species worldwide in seven families;
the Mexican species are about 40\% endemic
and 85\% belong to the Curculionidae
\cite{morrone2014curculionoidea}. Mexico has
about 14\,500 described and an estimated
23\,750 species of Lepidoptera, close to 10\%
of the world total
\cite{llorente2014lepidoptera}.
\item The marine molluscs are 2576 species on
the Pacific coast and 2067 in the Gulf of
Mexico and Caribbean, against 117\,358
described species worldwide
\cite{castillo2014moluscos}.
\item The fishes are the most species-rich
vertebrates of Mexico: 505 freshwater, 2224
marine and 563 estuarine and vicarious
species (the categories overlap)
\cite{espinosa2014peces}. ``Fishes'' are not
a natural group: tetrapods arose from within
the lobe-finned fishes, so the group is
paraphyletic \cite[Chapter
16]{hickman2002animal}.
\end{itemize}

\section{Chordates and vertebrates}
\label{section-chordates}

This section is written from the textbooks;
the surviving coursework has no material on
it.

\begin{definition}
\label{definition-chordate}
A \emph{chordate} is an animal that has, at
least at some embryonic stage, a
\emph{notochord}, a single dorsal tubular
nerve cord, pharyngeal pouches and a postanal
tail.
\end{definition}

The notochord is a rod of turgid, closely
packed, vacuolated cells inside fibrous and
elastic sheaths; it stiffens the body axis
and is replaced by the vertebral column in
most vertebrates, surviving in mammals only
in the nuclei pulposi of the intervertebral
discs \cite[Chapter 15]{hickman2002animal}.
Two subphyla lack a head and vertebrae:
Urochordata (tunicates, mostly sessile adults
with a swimming tadpole larva) and
Cephalochordata (lancelets such as
amphioxus). Vertebrates add a head with a
brain, a living endoskeleton that grows with
the animal, a pharynx perforated by gill
slits, and paired limbs \cite[Chapter
15]{hickman2002animal}.

\textbf{Fishes.} Hagfishes and lampreys are
jawless, eel-like, with no paired fins and a
notochord that persists all their lives. All
other vertebrates have jaws. Chondrichthyes
(sharks, rays, chimaeras) have a
cartilaginous skeleton, paired fins and
placoid scales. Bony fishes are the
ray-finned Actinopterygii, almost all of them
teleosts (about 23\,600 species), and the
lobe-finned Sarcopterygii, today only
lungfishes and the coelacanth, from one
lineage of which the tetrapods arose. Most
bony fishes regulate buoyancy with a
gas-filled swim bladder, and their gills
exchange gases by countercurrent flow
\cite[Chapter 16]{hickman2002animal}.

\textbf{Amphibians.} Ectothermic tetrapods
with moist glandular skin that breathe by
lungs, gills or skin. The living orders are
caecilians (Gymnophiona, limbless),
salamanders (Caudata, with a tail) and frogs
and toads (Anura, jumpers without a tail).
Most have an aquatic larva that
metamorphoses; some salamanders, like the
axolotl, stay permanently gilled
\cite[Chapter 17]{hickman2002animal}.

\textbf{Reptiles.} The \emph{amniotic egg},
with three extraembryonic membranes inside a
shell, freed the amniotes from water;
reptiles also have dry scaly skin and
internal fertilization. Early amniotes split
into the synapsids, which led to mammals, and
the reptilian lineages leading to turtles
(Testudines), lizards and snakes (Squamata,
95\% of living reptiles), the tuatara
(\emph{Sphenodon}) and the archosaurs:
crocodilians, dinosaurs and birds. As a
class, Reptilia is paraphyletic \cite[Chapter
18]{hickman2002animal}.

\textbf{Birds.} More than 9600 living species
of feathered, endothermic, egg-laying
vertebrates with forelimbs modified as wings.
Birds descend from theropod dinosaurs;
\emph{Archaeopteryx}, from the Jurassic, is
almost a feathered theropod. Flight is served
by weight reduction (feathers, hollow bones,
a horny bill instead of teeth) and by power:
a high metabolic rate and air sacs that pass
air through the lungs during both inspiration
and expiration \cite[Chapter
19]{hickman2002animal}.

\textbf{Mammals.} Endothermic vertebrates
with hair, mammary glands, a skin rich in
sweat, scent and sebaceous glands, and teeth
differentiated into incisors, canines,
premolars and molars, replaced once.
Monotremes lay eggs; marsupials give birth to
underdeveloped young that grow in a pouch;
eutherians nourish the embryo through a
placenta for a long gestation. Whales,
dolphins and porpoises (Cetacea, 78 species
in Hickman's count) have the forelimbs
modified into flippers, no hind limbs,
nostrils as one or two blowholes on top of
the head, teeth (when present) all alike and
without enamel, hair reduced to a few
bristles on the snout, and no external ear.
Baleen whales, the largest mammals, have a
gestation of only about twelve months
\cite[Chapter 20]{hickman2002animal}.

\bibliography{my}
\bibliographystyle{amsalpha}

\end{document}
